\documentclass[11pt,a4paper]{article}
\usepackage[T1]{fontenc}
\usepackage[utf8]{inputenc}
\usepackage[margin=25mm]{geometry}
\usepackage{amsmath,amssymb,amsthm,booktabs}
\usepackage[hidelinks]{hyperref}
\newtheorem{proposition}{Proposition}
\newcommand{\norm}[1]{\left\lVert#1\right\rVert}
\title{Gradient Descent and Heavy-Ball Momentum\\
\large Module 1: Spectral Baselines and Rotation Correspondence}
\author{CHENG Zixuan}
\date{16 October 2026}
\begin{document}
\maketitle

\begin{abstract}
We establish spectral baselines for gradient descent (GD) and heavy-ball
momentum on positive-definite quadratics. Both methods commute with orthogonal
coordinate changes when their initial states and scalar parameters correspond.
We derive GD's objective bound, heavy-ball's strict stability region, and the
optimal constant parameters for worst-mode asymptotic contraction. A reproducible
two-dimensional study checks 56 trajectories, including stability boundaries
and unstable controls, using an independent modal recurrence. These are module
verification results, rather than a tuned comparison of the group's methods.
\end{abstract}

\section{Problem and implementation convention}
Consider
\[
 f(w)=\tfrac12w^THw,\quad H=QDQ^T,\quad Q^TQ=I,\quad
 D=\operatorname{diag}(\lambda_1,\ldots,\lambda_p)\succ0.
\]
The minimizer is zero. Write $\mu=\min_i\lambda_i$, $L=\max_i\lambda_i$,
$\kappa=L/\mu$, $z=Q^Tw$, and $\phi(z)=\tfrac12z^TDz$.
The gradient identity is $\nabla f(Qz)=QDz$.
GD uses $w_{t+1}=w_t-\eta Hw_t$. Our heavy-ball convention is
\begin{equation}
 v_{t+1}=\beta v_t+Hw_t,\qquad
 w_{t+1}=w_t-\eta v_{t+1},\qquad v_0=0,
 \label{eq:hb}
\end{equation}
where $\eta>0$ and $0\le\beta<1$. Velocity stores accumulated gradients;
there is no $(1-\beta)$ normalization or bias correction. Eliminating it gives
\begin{equation}
 w_{t+1}=(1+\beta)w_t-\eta Hw_t-\beta w_{t-1},\qquad w_{-1}=w_0.
 \label{eq:position}
\end{equation}
The initial condition is part of the method, particularly when interpreting
transient amplification.

\section{Orthogonal correspondence}
\begin{proposition}
Run either GD or heavy-ball on $H=QDQ^T$ and on $D$, with identical scalar
parameters and corresponding initial points $w_0=Qz_0$. For heavy-ball, initialize
the velocities correspondingly, $v_0=Qu_0$ (both are zero here). Then in exact
arithmetic, for every step, $w_t=Qz_t$ and $f(w_t)=\phi(z_t)$.
\end{proposition}
\begin{proof}
For GD, $w_t=Qz_t$ implies
$w_{t+1}=Qz_t-\eta QDz_t=Q(z_t-\eta Dz_t)=Qz_{t+1}$.
For heavy-ball, if also $v_t=Qu_t$, then
$v_{t+1}=Q(\beta u_t+Dz_t)=Qu_{t+1}$ and
$w_{t+1}=Q(z_t-\eta u_{t+1})=Qz_{t+1}$.
Induction and $Q^THQ=D$ prove the statements.
\end{proof}
Gradient norms also correspond, since $\norm{QDz_t}_2=\norm{Dz_t}_2$.
The proof permits repeated eigenvalues and does not assume convergence. It
applies to orthogonal transformations and shared scalar parameters; it makes
no claim about coordinatewise adaptive methods. Finite arithmetic can perturb
trajectories, especially after unstable modes have amplified rounding errors.

\section{GD contraction and optimal constant step}
In eigen-coordinates,
\[
 z_{t,i}=(1-\eta\lambda_i)^t z_{0,i},\qquad
 \rho_{\rm GD}=\max_i|1-\eta\lambda_i|.
\]
Consequently
\begin{equation}
 f(w_t)=\tfrac12\sum_i\lambda_i(1-\eta\lambda_i)^{2t}z_{0,i}^2
 \le\rho_{\rm GD}^{2t} f(w_0).
 \label{eq:gd-bound}
\end{equation}
Convergence for every initial point holds exactly when $0<\eta<2/L$.
At $\eta=2/L$, the largest-eigenvalue mode alternates without decay; above
this boundary it grows if excited. An exceptional initial point without that
mode does not establish stability for all initial points.

The maximum of $|1-\eta\lambda|$ over the spectrum is attained at $\mu$ or $L$.
Balancing the endpoint values gives the minimax constant step and state factor
\[
 \eta^*_{\rm GD}=\frac2{L+\mu},\qquad
 \rho^*_{\rm GD}=\frac{L-\mu}{L+\mu}=\frac{\kappa-1}{\kappa+1}.
\]
The objective factor is its square. The common step $1/L$ instead gives
$\rho_{\rm GD}=1-1/\kappa$.

\section{Heavy-ball stability and minimax asymptotic factor}
Equation~\eqref{eq:position} yields, for each eigenvalue $\lambda$,
\[
 z_{t+1}=(1+\beta-\eta\lambda)z_t-\beta z_{t-1},\qquad
 p_\lambda(r)=r^2-(1+\beta-\eta\lambda)r+\beta.
\]
For $r^2-ar+\beta$, strict Schur stability is equivalent to
$1-a+\beta>0$, $1+a+\beta>0$, and $1-\beta>0$.
Substitution gives the necessary and sufficient condition, within our parameter
range, for convergence for all initial states:
\begin{equation}
 0\le\beta<1,\qquad 0<\eta<\frac{2(1+\beta)}L.
 \label{eq:stability}
\end{equation}
At equality there is a root $-1$ in the $L$ mode, hence no strict stability.
Our code checks the strict parameter inequality, since rounding can make a
computed boundary root modulus slightly smaller than one.

Let $\rho_{\rm HB}$ be the maximum root modulus over all modes. The optimal
constant parameters for this asymptotic criterion are
\begin{equation}
 \eta^*_{\rm HB}=\frac4{(\sqrt L+\sqrt\mu)^2},\quad
 \beta^*_{\rm HB}=\left(\frac{\sqrt L-\sqrt\mu}{\sqrt L+\sqrt\mu}\right)^2,
 \quad \rho^*_{\rm HB}=\frac{\sqrt\kappa-1}{\sqrt\kappa+1}.
 \label{eq:optimal-hb}
\end{equation}
Here is a direct optimality argument. Suppose all mode roots have modulus at
most $s<1$, with $s>0$. Their product gives $\beta\le s^2$. Evaluating the
polynomial at $s$ and $-s$ gives
$|1+\beta-\eta\lambda|\le s+\beta/s$.
Using the endpoint modes,
\[
 \eta\mu\ge(1-s)(1-\beta/s),\qquad
 \eta L\le(1+s)(1+\beta/s).
\]
Because $0\le\beta/s\le s<1$, it follows that
\[
 \kappa\le\frac{(1+s)(1+\beta/s)}{(1-s)(1-\beta/s)}
 \le\left(\frac{1+s}{1-s}\right)^2.
\]
Thus $s\ge(\sqrt\kappa-1)/(\sqrt\kappa+1)$. Parameters
\eqref{eq:optimal-hb} attain this bound: with $s=\rho^*_{\rm HB}$,
$\beta=s^2$ and the coefficients $1+\beta-\eta\lambda$ range from $2s$ to
$-2s$, so every root has modulus $s$. When $\mu=L$, both optimal methods reduce
to the exact one-step update with $\eta=1/L$, $\beta=0$.

This is an asymptotic state criterion, not the GD objective bound
\eqref{eq:gd-bound}. At the optimal endpoints the roots repeat and solutions
have the form $(A+Bt)r^t$. For the $L$ mode and zero initial velocity,
\begin{equation}
 z_t=[1+(1+\rho)t](-\rho)^t z_0,\qquad \rho=\rho^*_{\rm HB}.
 \label{eq:transient}
\end{equation}
Loss can increase initially despite strict stability. For $0<\rho<1$, the
repeated-root worst-case state envelope is $O((t+1)\rho^t)$ and the objective
envelope is $O((t+1)^2\rho^{2t})$, with problem-dependent constants. A uniform
claim $f(w_t)\le\rho^{2t}f(w_0)$ would be false.

\section{A hand-checkable two-dimensional example}
Let $D=\operatorname{diag}(0.01,1)$, $z_0=(1,1)^T$ and let $Q$ rotate by
$45^\circ$. Then
\[
 H=\begin{pmatrix}0.505&-0.495\\-0.495&0.505\end{pmatrix},\quad
 w_0=(0,\sqrt2)^T,\quad f(w_0)=0.505.
\]
With $\eta=1$, both methods have $z_1=(0.99,0)^T$, $w_1=Qz_1$, and loss
$0.0049005$. With heavy-ball $\beta=0.5$, the next point is
$z_2=(0.9751,-0.5)^T$ and the loss is $0.12975410005$. The method is stable
but nonmonotone. Dedicated tests check these values in rotated coordinates.

\section{Observed module verification}
The run of 16 October 2026 used Python 3.11.17 and NumPy 2.4.6, seed 5020,
spectrum $(0.01,1)$, angles $0,30,45,90$ degrees, 300 updates, target relative
loss $10^{-6}$, and explosive threshold $10^{12}$. A single seeded initial
direction is normalized to unit loss and shared across rotations via $w_0=Qz_0$.
Seven parameter settings at four angles give 28 pairs and 56 saved trajectories.
Table~\ref{tab:results} gives diagonal-reference results; rotated runs have
matching statuses and attainment steps.

\begin{table}[htbp]
\centering
\caption{Fixed theory checks, without parameter search. A dash means the
target was not reached within the recorded run.}
\label{tab:results}
\small
\begin{tabular}{lrrrl}
\toprule
Setting & Calls & Final relative loss & First target & Status\\
\midrule
GD, $\eta=1/L$ & 300 & $3.30005\,10^{-5}$ & -- & finite cap\\
GD, optimal & 300 & $6.14176\,10^{-6}$ & -- & finite cap\\
GD, boundary & 300 & $0.986279$ & -- & finite cap\\
GD, outside boundary & 146 & $1.20410\,10^{12}$ & -- & explosive\\
HB, optimal & 300 & $1.51030\,10^{-47}$ & 58 & finite cap\\
HB, boundary & 300 & $25.1526$ & -- & finite cap\\
HB, outside boundary & 45 & $1.11170\,10^{12}$ & -- & explosive\\
\bottomrule
\end{tabular}
\end{table}

The optimal GD parameters are $(\eta,\rho)=(1.980198,0.980198)$; optimal HB
uses $(\eta,\beta,\rho)=(3.305785,0.669421,0.818182)$.
Boundary controls use $2/L$ and $2(1+\beta)/L$; unstable controls multiply
these boundary steps by $1.05$. Stable paired paths have maximum discrepancy
\[
 \max_t\norm{w_t-Qz_t}/\max(1,\norm{z_0})=6.240\,10^{-15}.
\]
Including growing unstable controls raises this metric to $6.849\,10^{-9}$;
the trajectories then have magnitudes around one million. Checks use mixed
componentwise absolute and relative tolerances $10^{-9}$ after fixed initial
scaling, rather than a universal $10^{-9}$ absolute-error requirement.
The largest loss discrepancy scaled by $\max(1,|\phi(z_t)|)$ is
$3.209\,10^{-13}$.

An independent checker regenerated the seeded initial point and all schedules,
then reconstructed every saved trajectory through modal position recurrences.
It imports neither the production runner nor the update and spectral helpers.
It checked losses, gradient norms, stopping, target attainment, CSV summaries,
operation counts and source hashes. Its maximum path reconstruction discrepancy,
scaled at each step by the larger of initial and current state norms and one,
was $1.204\,10^{-13}$; all 56 trajectories passed.

The public runner records step zero before updates. Reaching the target does
not stop a run. A finite-cap status means the budget finished without numerical
failure or an explosive event, not that convergence was achieved. Unattained
targets remain censored by the budget; explosive trajectories retain their
triggering step and loss. Each update uses one gradient call; monitoring losses
are additional evaluation work. Neither optimizer needs an eigenbasis to run,
but selecting the theoretical optimal parameters uses spectral endpoints.

\section{Reproduction and integration}
From the repository root in the project environment, run
\begin{verbatim}
python -m unittest discover -s tests -v
python -m coordinate_geometry.part1_experiments \
  --config configs/part1.json --output outputs/part1
python -m coordinate_geometry.verify_part1 \
  --output outputs/part1
\end{verbatim}
The first experiment command can take \texttt{--no-plots}. Artifacts under
\path{outputs/part1/} include:
\begin{flushleft}
\path{summary.csv}, \path{trajectories.json}, \path{run_config.json},\\
\path{summary.json}, \path{artifact_verification.json}, and\\
\path{part1_validation.png}.
\end{flushleft}
The JSON summary records environment versions, a UTC generation timestamp,
and source/\allowbreak configuration SHA256 hashes. These identify the local source
snapshot even in an extracted folder without Git metadata. This standalone
report source has no external figure dependency.

The integration identifiers are \texttt{gd} and \texttt{heavy\_ball}, with
\texttt{learning\_rate} and \texttt{beta} in the unchanged shared contract.
The completed baselines explain why rotating a fixed spectrum cannot change
these methods' idealized paired trajectories. Their optimal-parameter example
does not establish universal finite-budget or wall-clock superiority of
heavy-ball. The group must still freeze its pilot protocol, perform team
cross-review, and compare all methods under the agreed tuning and statistical
rules before drawing main-study conclusions.
\end{document}
